\Needspace{12\baselineskip}
\section{西安算例的补充结果}
\label{sec:sup:xian}
本节补充正文第~\ref{sec:xian}~节的西安算例：先给出初值对照与阈值扫描，再在附加的医疗需求假设下给出阈值的一个充分条件（第~\ref{sec:sup:medical-link}~节，正文第~\ref{sec:model}~节开头引用），最后说明设计阈值 $\eta=0.002N$ 的量级（第~\ref{sec:sup:threshold-scale}~节）。西安参数见正文表~\ref{tab:xian_initial_fit}，计算使用 $I_0$ 的未取整值。

\subsection{初值对照与阈值扫描}
\label{sec:sup:xian:settings}
拟合得到的 $I_0$ 远小于 $1$。表~\ref{tab:sup:initial}把它与固定 $I_0=1$ 的对照放在一起：后者使观测窗口内的累计病例高出两个数量级以上，与观测数据不符（正文第~\ref{sec:xian}~节）。
\begin{table}[htbp]
\centering\small
\setlength{\tabcolsep}{3.5pt}
\renewcommand{\arraystretch}{1.25}
\begin{threeparttable}
\caption{不同初值下观测窗口内的模型结果}
\label{tab:sup:initial}
\begin{tabular*}{\linewidth}{@{\extracolsep{\fill}}lS[table-format=1.6]S[table-format=7.0]S[table-format=5.2]S[table-format=5.2]@{}}
\toprule
数据或初值设定 & {$I_0$（人）} & {\shortstack{40日累计新增病例\\（例）}} & {\shortstack{社区日新增峰值\\（例）}} & {\shortstack{隔离日新增峰值\\（例）}}\\
\midrule
观测数据 & {---} & 2050 & 60 & 141\\
最小二乘拟合 & 0.001007 & 2099 & 53.88 & 115.93\\
固定 $I_0=1$ & 1 & 1083553 & 36873.64 & 63966.50\\
\bottomrule
\end{tabular*}
\begin{tablenotes}[flushleft]\footnotesize
\item 注：观测窗口为 2021 年 12 月 9 日至 2022 年 1 月 17 日，共 40 天；拟合行使用未取整的 $I_0$。日新增峰值是病例流量的峰值，不同于现存感染者的峰值。TDINN 拟合轨迹的总累计新增感染在 $t=40$ d 为 $2098.9747$，在清零时间 $45.2874$ d 为 $2105.6330$；收紧积分容差并把最大步长减半后，后者的变化为 $3.61\times10^{-8}$ 人。
\end{tablenotes}
\end{threeparttable}
\end{table}

表~\ref{tab:sup:eta}与图~\ref{fig:xian_eta_sens}给出 11 个阈值比例下仅隔离控制的结果，覆盖设计阈值 $\eta/N=0.002$ 的两侧。阈值从 $0.01N$ 降到 $0.0001N$ 时，$q_{\max}$ 只在 $0.838$ 与 $0.847$ 之间变化，控制时长与成本近似按阈值的反比增加，总累计新增感染则缓慢减少；正文第~\ref{sec:xian:threshold-scale}~节由此解释低阈值的代价。
\begin{table}[htbp]
\centering\small
\setlength{\tabcolsep}{3.5pt}
\renewcommand{\arraystretch}{1.25}
\begin{threeparttable}
\caption{西安参数下不同阈值的仅隔离控制}
\label{tab:sup:eta}
\begin{tabular*}{\linewidth}{@{\extracolsep{\fill}}S[table-format=1.4]S[table-format=1.4]S[table-format=4.2]S[table-format=4.2]S[table-format=3.4]S[table-format=7.0]@{}}
\toprule
{$\eta/N$} & {$q_{\max}$} & {$\Delta t$（d）} & {$t_{\rm end}$（d）} & {$J$（d）} & {\shortstack{总累计新增感染\\（人）}}\\
\midrule
0.0001 & 0.8470 & 1710.66 & 2209.55 & 826.4066 & 2155649\\
0.0002 & 0.8469 & 855.09 & 1243.15 & 412.9141 & 2182123\\
0.0005 & 0.8466 & 341.75 & 620.77 & 164.8184 & 2234666\\
0.0010 & 0.8462 & 170.63 & 389.32 & 82.1192 & 2293939\\
0.0015 & 0.8458 & 113.59 & 304.01 & 54.5524 & 2339486\\
0.0020 & 0.8454 & 85.07 & 258.10 & 40.7686 & 2377943\\
0.0030 & 0.8445 & 56.55 & 208.39 & 26.9842 & 2442608\\
0.0040 & 0.8436 & 42.29 & 181.18 & 20.0913 & 2497306\\
0.0050 & 0.8428 & 33.73 & 163.65 & 15.9550 & 2545663\\
0.0075 & 0.8405 & 22.32 & 138.10 & 10.4383 & 2649240\\
0.0100 & 0.8382 & 16.61 & 123.89 & 7.6781 & 2737382\\
\bottomrule
\end{tabular*}
\begin{tablenotes}[flushleft]\footnotesize
\item 注：$q_{\max}=q_c(t_1)$；$\eta/N$ 为比例，$0.002$ 即 $0.2\%$；总累计新增感染计算至清零时间。
\end{tablenotes}
\end{threeparttable}
\end{table}

\begin{figure}[htbp]
    \centering
    \includegraphics[width=0.7\textwidth]{xian_eta_sensitivity.pdf}
    \caption{西安参数下仅隔离控制的指标随 $\eta/N$ 的变化：(a) 启动时的隔离比例；(b) 成本 $J$；(c) 累计新增感染；(d) 控制时长与清零时间。水平线为 TDINN 参照值；(d) 中 TDINN 的控制时长与清零时间同为 $45.29$ d。}
    \label{fig:xian_eta_sens}
\end{figure}

图~\ref{fig:xian_heatmaps}把扫描推广到 $(\eta,c_0)$ 平面。在图示的 $c_0$ 范围内，控制时长、清零时间与成本的等值线接近竖直，主要由阈值决定；正文第~\ref{sec:xian:threshold-scale}~节在 TDINN 峰值处的控制时长即取自子图 (a)。
\begin{figure}[htbp]
    \centering
    \includegraphics[width=0.8\textwidth]{xian_heatmaps.pdf}
    \caption{西安参数下仅隔离控制在 $(\eta,c_0)$ 平面上的指标：(a) 控制时长；(b) 清零时间；(c) 累计新增感染；(d) 成本 $J$。两条竖直参照线分别为 TDINN 的社区感染峰值 $\eta=152.55$ 与设计阈值 $\eta=26326$；虚线为等值线。}
    \label{fig:xian_heatmaps}
\end{figure}

\FloatBarrier
\subsection{阈值与医疗需求}
\label{sec:sup:medical-link}
正文第~\ref{sec:model}~节开头指出，接触率不超过常规水平时，限制 $I$ 也就限制了进入 $I$ 和 $I_q$ 的全部新增感染。本节给出这一上界，并在附加的医疗需求假设下由容量导出阈值的一个充分条件。

两类感染仓室的新增流入之和为
\[
\frac{\beta c(t)(1-q(t))S(t)I(t)}{N}
+\frac{\beta c(t)q(t)S(t)I(t)}{N}
=\frac{\beta c(t)S(t)I(t)}{N}.
\]
若全过程满足 $0\le c(t)\le c_0$ 和 $I(t)\le\eta$，由 $S(t)\le S_0$，得到
\[
0\le\frac{\beta c(t)S(t)I(t)}{N}
\le\frac{\beta c_0S_0}{N}\eta.
\]
这一上界包含进入 $I$ 和 $I_q$ 的全部新增感染。

设非负可测函数 $h(a)$ 表示一例新感染者在感染后第 $a$ 天所需的平均 ICU 床位数，并假定它对两类病例、不同感染时点及所比较的控制共同适用。令
\[
\int_0^\infty h(a)\,da=pL<\infty,
\]
其中 $p\in(0,1]$ 为需要 ICU 的概率，$L>0$ 为需要 ICU 者累计占用时长的均值；$h$ 可包含感染至入院的延迟。由既往流入与停留过程计算床位占用的方法可参见文献\cite{Baas2021Occupancy}。

初始时刻之前已经感染者此后的平均需求记为 $D_0(t)$，并假定各策略均满足 $0\le D_0(t)\le D_{\mathrm{init}}<\infty$。该项包括初始病例以后才发生的入院和后续占用，也包括已从传播模型移出但仍需救治的病例。不受容量拒收限制的平均需求为
\[
D(t)=D_0(t)+\int_0^t h(t-u)
\frac{\beta c(u)S(u)I(u)}{N}\,du.
\]
利用非负性及总流入上界，有
\[
\begin{aligned}
D(t)
&\le D_{\mathrm{init}}+
\frac{\beta c_0S_0}{N}\eta\int_0^t h(a)\,da\\
&\le D_{\mathrm{init}}+pL\frac{\beta c_0S_0}{N}\eta,
\qquad t\ge0.
\end{aligned}
\]
设 $C>D_{\mathrm{init}}$ 为全过程固定的本病可用容量，已扣除其他病种占用和安全预留，本病初始病例的未来负担由 $D_{\mathrm{init}}$ 单独计入。于是
\[
\eta\le\frac{C-D_{\mathrm{init}}}{pL\,\beta c_0S_0/N}
\]
是平均需求满足 $D(t)\le C$ 的充分条件。该条件对 $\mathcal A(\eta)$ 中的全部控制同时成立，因而可能偏保守；它约束的是平均需求，不给出随机超载的概率。按此选取的阈值还须满足触发条件\eqref{eq:model:trigger}与能力约束下的可行条件（正文推论~\ref{cor:joint:feasibility}）。

\subsection{设计阈值的量级}
\label{sec:sup:threshold-scale}
正文西安算例取 $N=13{,}163{,}000$\cite{XianStat2021}、$\eta=0.002N=26{,}326$。这一阈值来自下面的量级估计，它与第~\ref{sec:sup:medical-link}~节的充分条件是两种不同的依据。

陈胤孜等\cite{Chen2021ICUForecast}基于既往资料预测，2021 年全国每10万人综合 ICU 床位约为 $4.37$ 张；这一数值为全国预测量，而非西安同期可用容量的实测值。Angulo 等\cite[补充材料 S4.1]{Angulo2021}对早期 COVID-19 采用需要重症监护的感染者比例名义值 $0.60\times0.15\times0.28=0.0252$。将前者作为人均容量参考、不扣除既有占用，并将后者移作参考系数 $\rho_{\mathrm{ICU}}$，按比例情景
\[
\eta=\frac{C}{\rho_{\mathrm{ICU}}},\qquad
\theta=\frac{C/N}{\rho_{\mathrm{ICU}}}
\]
得到
\[
\theta=\frac{4.37/10^5}{0.0252}\approx1.73\times10^{-3},
\qquad \eta\approx2.28\times10^4.
\]
正文在这一量级附近取 $\theta=0.002$ 作为设计情景。按相同的参考系数反算，它对应每10万人 $5.04$ 张、全市约 $663$ 张的名义容量。

这一比例移植只给出阈值的量级，不是当地医疗容量的临床校准。以确诊病例或全部感染者为分母的 ICU 入院风险可能不同\cite{Pung2023Undetected}；参考系数 $\rho_{\mathrm{ICU}}$ 也没有依据共同的需求函数和初始负担估计，不等同于第~\ref{sec:sup:medical-link}~节中的 $pL\beta c_0S_0/N$。
